\documentclass[12pt]{article}

\usepackage{amsmath, amssymb, amsthm}
\usepackage{enumitem}
\usepackage{geometry}
\usepackage{fancyhdr}   % headers/footers
\usepackage{graphicx}
\usepackage{hyperref}

% --- Page Setup ---
\geometry{margin=1in}
\pagestyle{fancy}
\fancyhf{}
\rhead{MATH 421 -- 003, Fall 2025}
\lhead{Homework 4}
\cfoot{\thepage}
\renewcommand{\headrulewidth}{0.4pt}

% --- Title info ---
\title{Homework 4}
\author{Alejandro De La Torre \\ MATH 421 -- The Theory of Single Variable Calculus \\ Section 003 \\ Instructor: Grace Work}
\date{October 2025}

% --- Theorem environments (handy if you need them) ---
\newtheorem{theorem}{Theorem}
\newtheorem{lemma}{Lemma}
\newtheorem{proposition}{Proposition}
\newtheorem{corollary}{Corollary}

\begin{document}
\maketitle

\section*{Collaboration Statement}
Work compared with Griffin Hailman, with whom brief advice and discussion were exchanged on some problems. Mainly referenced the textbook \textit{Elementary Analysis: The Theory of Calculus} by Ross and my own class notes. 

% ---------------- Problem 1 ----------------
\section*{Problem 1 (Ross Chapter 1, 4.7)}
Let $S$ and $T$ be nonempty bounded subsets of $\mathbb{R}$.

\medskip
\noindent\textbf{(a)} If $S\subseteq T$, then
\[
\inf T \;\le\; \inf S \;\le\; \sup S \;\le\; \sup T .
\]

\begin{proof}
Assume $S\subseteq T$. We verify the three inequalities separately.
\begin{enumerate}[label=\roman*)]
  \item \emph{$\sup S\le \sup T$.}  Every element of $S$ lies in $T$, so any
  upper bound of $T$ is an upper bound of $S$. In particular,
  $\sup T$ is an upper bound of $S$. By minimality of the least upper bound,
  $\sup S\le \sup T$.

  \item \emph{$\inf T\le \inf S$.}  Again $S\subseteq T$, hence any lower bound
  of $T$ is a lower bound of $S$. In particular, $\inf T$ is a lower
  bound of $S$. By maximality of the greatest lower bound, $\inf T\le \inf S$.

  \item \emph{$\inf S\le \sup S$.}  Because $S$ is nonempty and bounded,
  there exists $s\in S$ and every $s\in S$ satisfies
  $\inf S\le s\le \sup S$. Hence $\inf S\le \sup S$.
\end{enumerate}
Putting the three statements together yields
$\inf T \le \inf S \le \sup S \le \sup T$.
\end{proof}

\medskip
\noindent\textbf{(b)} Without assuming $S\subseteq T$,
\[
\sup(S\cup T)=\max\{\sup S,\sup T\}.
\]

\begin{proof}
Let $\alpha:=\sup S$ and $\beta:=\sup T$, and set $M:=\max\{\alpha,\beta\}$.

\emph{Step 1: $M$ is an upper bound for $S\cup T$.}
If $x\in S\cup T$, then either $x\in S$ or $x\in T$.
In the first case $x\le \alpha\le M$; in the second $x\le \beta\le M$.
Hence $x\le M$ for all $x\in S\cup T$.

\emph{Step 2: $M$ is the least such upper bound.}
Let $u$ be any upper bound for $S\cup T$.
Then $u$ is an upper bound for $S$ and for $T$, so
$\alpha=\sup S\le u$ and $\beta=\sup T\le u$. Consequently
$M=\max\{\alpha,\beta\}\le u$.

By Steps~1–2, $M$ is an upper bound for $S\cup T$ that is no larger
than any other upper bound; by definition, $M=\sup(S\cup T)$.
\end{proof}

% ---------------- Problem 2 ----------------
\section*{Problem 2 (Ross Chapter 1, 4.8)}
Let $S$ and $T$ be nonempty subsets of $\mathbb{R}$ with the property
\[
(\ast)\qquad s\le t \quad\text{for all } s\in S \text{ and } t\in T .
\]

\begin{enumerate}[label=\textbf{(\alph*)}]

\item \textbf{Observation. $S$ is bounded above and $T$ is bounded below.}

\emph{Reason.}
Pick any $t_0\in T$ (possible because $T\neq\varnothing$).  By $(\ast)$, for every
$s\in S$ we have $s\le t_0$, so $t_0$ is an upper bound of $S$ and hence $S$
is bounded above.

Likewise, pick any $s_0\in S$.  Then for every $t\in T$ we have $s_0\le t$, so
$s_0$ is a lower bound of $T$ and hence $T$ is bounded below. \hfill $\square$

\item \textbf{} $\displaystyle \sup S \le \inf T$.

\begin{proof}
By part (a), $\sup S$ and $\inf T$ exist (since $S\neq\varnothing$ is bounded
above and $T\neq\varnothing$ is bounded below).
Fix $t\in T$.  Property $(\ast)$ implies $s\le t$ for all $s\in S$, i.e., $t$ is an
upper bound for $S$.  Hence $\sup S\le t$ for every $t\in T$.
Therefore $\sup S$ is a lower bound for $T$, and by maximality of $\inf T$,
\(\sup S \le \inf T\).
\end{proof}

\item \textbf{Example with $S\cap T\neq\varnothing$.}

Take
\[
S=[0,1], \qquad T=[1,2].
\]
Then for any $s\in S$ we have $s\le 1\le t$ for every $t\in T$, so $(\ast)$ holds.
Moreover, $S\cap T=\{1\}\neq\varnothing$.

\item \textbf{Example with $\sup S=\inf T$ and $S\cap T=\varnothing$.}

Take
\[
S=(0,1), \qquad T=[1,2] \quad\text{(or }T=(1,2]\text{)}.
\]
For any $s\in(0,1)$ and any $t\in T$ we have $s<1\le t$, so $(\ast)$ holds.
Here $\sup S=1=\inf T$, while $S\cap T=\varnothing$ because $1\notin S$ (and if
$T=(1,2]$ then $1\notin T$ either).
\end{enumerate}

% ---------------- Problem 3 ----------------
\section*{Problem 3 (Ross Chapter 1, 4.12)}
Let $\mathbb{I}$ be the set of real numbers that are not rational; elements of $\mathbb{I}$ are called \emph{irrational numbers}.  
Prove: if $a<b$, then there exists $x\in\mathbb{I}$ such that $a<x<b$.

\begin{lemma}[A rational in any open interval]
\label{lem:densityQ}
If $u<v$ are real numbers, then there exists $r\in\mathbb{Q}$ with $u<r<v$.
\end{lemma}

\begin{proof}
By the Archimedean property, choose $n\in\mathbb{N}$ with $n(v-u)>1$.  
There is an integer $m$ with $nu<m<nv$ (e.g., take $m=\lfloor nv\rfloor$; then $m<nv$ and $m>nv-1>nu$).  
Set $r=\dfrac{m}{n}\in\mathbb{Q}$. Then
\[
u<\frac{nu}{n}<\frac{m}{n}<\frac{nv}{n}<v,
\]
so $u<r<v$.
\end{proof}

\begin{lemma}\label{lem:shift}
If $r\in\mathbb{Q}$, then $r+\sqrt2\in\mathbb{I}$.
\end{lemma}

\begin{proof}
Suppose, toward a contradiction, that $r+\sqrt2\in\mathbb{Q}$.  
Since $\mathbb{Q}$ is closed under subtraction (is a field where the difference of rationals is rational?), $(r+\sqrt2)-r=\sqrt2$ would be rational, which contradicts the well–known fact that $\sqrt2$ is irrational.  
Hence $r+\sqrt2$ is irrational; i.e., $r+\sqrt2\in\mathbb{I}$.
\end{proof}

For any $a<b$ in $\mathbb{R}$, there exists $x\in\mathbb{I}$ with $a<x<b$.

\begin{proof}
Apply Lemma~\ref{lem:densityQ} to the interval $(a-\sqrt2,\; b-\sqrt2)$ to obtain
a rational $r\in\mathbb{Q}$ such that
\[
a-\sqrt2 \;<\; r \;<\; b-\sqrt2.
\]
Let $x:=r+\sqrt2$.  Adding $\sqrt2$ to the inequalities gives
\[
a \;<\; x \;<\; b.
\]
By Lemma~\ref{lem:shift}, $x=r+\sqrt2$ is irrational; hence $x\in\mathbb{I}$, as required.
\end{proof}

\newpage

% ---------------- Problem 4 ----------------
\section*{Problem 4 (Ross Chapter 1, 4.14(a))}

4.14 Let $A$ and $B$ be nonempty bounded subsets of $\mathbb{R}$, and let 
$A + B$ be the set of all sums $a + b$ where $a \in A$ and $b \in B$.  

\begin{enumerate}[label=\textbf{(\alph*)}]
  \item Prove $\sup(A+B) = \sup A + \sup B$. 
  \textit{Hint:} To show $\sup A + \sup B \le \sup(A+B)$, show that for each $b \in B$, 
  $\sup(A+B) - b$ is an upper bound for $A$, hence $\sup A \le \sup(A+B) - b$. 
  Then show $\sup(A+B) - \sup A$ is an upper bound for $B$. 
\end{enumerate}

\begin{proof}
Since $A$ and $B$ are nonempty and bounded, $\sup A$ and $\sup B$ exist.  
Set $\alpha:=\sup A$ and $\beta:=\sup B$.

We prove the two inequalities separately.

\smallskip
\noindent\emph{(i) $\sup(A+B)\le \alpha+\beta$.}
For any $a\in A$ and $b\in B$ we have $a\le \alpha$ and $b\le \beta$, hence
\[
a+b\le \alpha+\beta.
\]
Thus $\alpha+\beta$ is an upper bound for $A+B$, and by minimality of the least upper bound,
\[
\sup(A+B)\le \alpha+\beta.
\]

\smallskip
\noindent\emph{(ii) $\sup(A+B)\ge \alpha+\beta$.}
Fix $b\in B$. For every $a\in A$ we have $a+b\le \sup(A+B)$, so
\[
a\le \sup(A+B)-b\qquad(a\in A).
\]
Hence $\sup(A+B)-b$ is an upper bound for $A$, and therefore
\[
\alpha=\sup A\le \sup(A+B)-b \quad\text{for every } b\in B,
\]
which implies $\alpha+b\le \sup(A+B)$ for every $b\in B$. Taking suprema over $b$ gives
\[
\sup_{b\in B}(\alpha+b)\le \sup(A+B).
\]
Since adding a constant preserves order, \(\sup_{b\in B}(\alpha+b)=\alpha+\sup B=\alpha+\beta\). Hence
\[
\alpha+\beta\le \sup(A+B).
\]

\smallskip
Combining (i) and (ii) yields \(\sup(A+B)=\alpha+\beta=\sup A+\sup B\).
\end{proof}

\vspace{3cm}

% ---------------- Problem 5 ----------------
\section*{Problem 5 (Ross Chapter 1, 4.15)}
Let $a,b\in\mathbb{R}$. Show: if $a\le b+\dfrac{1}{n}$ for all $n\in\mathbb{N}$, then $a\le b$. 
(Compare Exercise 3.8.)

\begin{proof}
We argue by contradiction. Suppose $a>b$. Set $\varepsilon:=a-b>0$.  
By the Archimedean property, there exists $n\in\mathbb{N}$ with
\[
\frac{1}{n}<\varepsilon=a-b.
\]
Then
\[
b+\frac{1}{n}\;<\;b+(a-b)\;=\;a,
\]
which contradicts the hypothesis $a\le b+\dfrac{1}{n}$ (true for every $n\in\mathbb{N}$).  
Therefore our assumption $a>b$ is impossible, and we conclude $a\le b$.
\end{proof}


% ---------------- Problem 6 ----------------
\section*{Problem 6 (Ross Chapter 1, 4.16)}
Show
\[
\sup\{\, r \in \mathbb{Q} : r < a \,\} \;=\; a
\quad\text{for each } a \in \mathbb{R}.
\]

\begin{proof}
Let 
\[
S := \{\, r \in \mathbb{Q} : r < a \,\}.
\]

\begin{enumerate}[label=\textbf{Step \arabic*.}]
  \item \emph{\(S\) is nonempty and bounded above.}  
  By the density of \(\mathbb{Q}\) in \(\mathbb{R}\), for any \(u<v\) there exists \(q\in\mathbb{Q}\) with \(u<q<v\).  
  Taking \(u=a-1\), \(v=a\), we get \(q<a\), so \(S\neq\varnothing\).  
  Since every \(r\in S\) satisfies \(r<a\), the real number \(a\) is an upper bound for \(S\).  
  Hence \(\sup S\) exists and \(\sup S \le a\).
  
  \item \emph{\(a\) is the least upper bound.}  
  Let \(\varepsilon>0\). By the density of \(\mathbb{Q}\), there exists \(q\in\mathbb{Q}\) with
  \[
  a-\varepsilon < q < a.
  \]
  Then \(q\in S\) and \(q>a-\varepsilon\). Thus \(a-\varepsilon\) is not an upper bound of \(S\).  
  Therefore \(\sup S \ge a-\varepsilon\).  
  Since this is true for every \(\varepsilon>0\), we conclude \(\sup S \ge a\).
\end{enumerate}

Combining Step~1 and Step~2 gives
\[
\sup S = a.
\]
\end{proof}

\end{document}